\section{Discussion and Limitations}
\label{sec:discussion}

The ablations recorded in run~E (Table~\ref{tab:ablations}) are consistent with the locked comparisons. At equal costs over the matched histories, the pooled-conformal demand quantiles had a mean cost of 0.3984 against 0.4075 for the raw-sales point model with residual quantiles. Pooled calibration moved q90 coverage from 0.8121 to 0.8832 against a nominal 0.90. The adaptive variant cost 0.4039 against 0.3984 for static calibration, so its drift correction did not pay off in this week. A hypothetical two-day carry-over of unsold stock lowered the cost from 0.3984 to 0.2579, which shows how much the cost level depends on the one-day shelf-life assumption. Giving the model the target-day weather and promotion values, which are unknown when the order is placed, changed the q50 WAPE only from 33.48\% to 33.40\%. A reduced-lag ablation was not run, because fitting a new model after the final week had been scored would fall outside the lock.

\begin{table}[t]
\centering\scriptsize
\caption{Ablations Recorded in Run~E, as Pairs of Values in the Order Named (Mean Order Cost at Equal Costs Unless the Scope Says Otherwise)}
\label{tab:ablations}
\begin{tabularx}{\linewidth}{@{}>{\raggedright\arraybackslash}Xl>{\raggedright\arraybackslash}p{0.34\linewidth}@{}}
\toprule
Recorded diagnostic & Values & Scope \\
\midrule
Raw-sales vs.\ demand-quantile & 0.4075 / 0.3984 & Equal-cost matched cells \\
Uncalibrated vs.\ pooled q90 coverage & 0.8121 / 0.8832 & Final donor proxy; nominal 0.90 \\
Direct vs.\ recovered-history TFT & 0.4999 / 0.4690 & Equal-cost deep subset \\
Short vs.\ full lag set & Not available & Not run after the lock \\
Static vs.\ adaptive calibration & 0.3984 / 0.4039 & Equal-cost sensitivity \\
One-day vs.\ two-day inventory & 0.3984 / 0.2579 & Hypothetical carry-over \\
Contract vs.\ oracle-weather WAPE & 33.48\% / 33.40\% & Target-day oracle; not deployable \\
\bottomrule
\end{tabularx}

\end{table}

The main limitations are the following.
\begin{itemize}
\item \emph{Selected proxy.} Donor-day sales equal demand only on complete days, and which days are complete depends on demand and stocking. The proxy does not identify demand on censored days, and the real-history analysis gives only a lower bound.
\item \emph{Simulated mechanisms.} Mechanisms A, B, and C are specified simulations, not counterfactual inventory paths. The transfer results measure sensitivity to these constructed histories, not transport to another retailer or stocking policy.
\item \emph{Different populations.} The neural policies cover the 600-series subset, and Kaplan--Meier scores only estimable rows. A comparison of all policies on a common population would be a new analysis after the final week was exposed and would have to be labeled as such.
\item \emph{Cost model.} Rankings depend on the locked cost ratios and on the one-day shelf life.
\item \emph{Deep models.} The TFT and TimesNet routes were trained with a small CPU and epoch budget, and the official baseline was reproduced only in part.
\item \emph{Final-week exposure.} The final week was opened once under the Phase~8 lock. Any later protocol evaluated on it is exploratory. A future evaluation on all 50,000 series would contain the 5,000 series already scored and should report the remaining series separately.
\item \emph{Intervals.} The intervals come from only 200 bootstrap replicates, so their tails are approximate; dates are not resampled separately, and no multiplicity adjustment is made.
\item \emph{Provenance.} No run hashed the notebook source it executed, and the upstream dataset revision and download date are unknown; the local file hashes are the operative identifiers. All runs executed on Linux with Python 3.11.15, two logical CPUs, and 8.42~GB of RAM, run~E without a GPU. The released outputs omit large intermediate files, so Phase~5 cannot be resumed from them without regenerating Phases~0--4.
\end{itemize}

Every table and figure is generated by a script from the saved output files of the three runs, without retraining any model.

\section{Conclusion}
\label{sec:conclusion}

A matched donor-proxy design makes stockout-aware forecasting and ordering comparisons on FreshRetailNet-50K checkable: models trained on differently censored histories are scored against the same outcomes, under a protocol fixed before the final week was read. Within this design, the value of calibrated demand quantiles depends on the cost ratio. They gave the lowest cost when shortage was more costly, while a raw-sales point model with residual quantiles was cheaper when excess stock was more costly. Evaluating the calibrated policy on an unseen mechanism raised its cost by at most 0.0472 normalized units, and TimesNet-recovered histories lowered TFT order costs within the deep subset. These findings hold for donor days, the simulated mechanisms, and the locked cost settings. Testing them on censored demand itself, or on series outside the evaluated sample, requires a new protocol.

\section*{Code and Data Availability}
\label{sec:availability}

The code, the saved run outputs, and the paper sources are available at \repourl. The repository holds the notebooks, the saved run outputs used here, and the scripts that regenerate the tables and figures. Runs~S, P, and~E are stored under \path{results/phase0_4/demo5000/}, \path{results/phase0_4/full50k/}, and \path{results/phase0_9_5k/}; in run~E, \path{manifests/final_test_lock.json} holds the locked protocol and \path{phase9_evaluation/final_order_metrics.csv} the scored population and row count of every policy. The repository does not redistribute the dataset, which is available from Hugging Face~\cite{dingdong2025datasetcard} under its CC BY 4.0 license. The upstream revision and download date of our copies were not recorded, so the local files are identified by their SHA-256 hashes below. The last line is the SHA-256 digest of the locked protocol content that run~E stored in its Phase~8 lock.
\begin{flushleft}\scriptsize
\texttt{train.parquet}\\
\texttt{6706832db892bbae4969c19d87e07975d2543d2ba7d7d4756360654785de5a3d}\\[2pt]
\texttt{eval.parquet}\\
\texttt{1b118840664280c6b88bffc84c80ee1f54c05d911e354b7599e5da10995e960e}\\[2pt]
\texttt{protocol lock digest}\\
\texttt{4ab12c4bddc06b4bb32ba8825353071b964ef5bfdd423bd4c876ac382180c1b6}
\end{flushleft}
